\epigraph{Fundamentación Teórica}{
La base teórica funciona como el armazón conceptual sobre el cual se erige todo
trabajo investigativo. Esta sección delimita las nociones esenciales, los
fundamentos teóricos existentes y los antecedentes modelares que respaldan la
indagación, brindando un sistema de referencia que guía la captación y
interpretación de los datos. Este epígrafe establece el marco conceptual necesario 
para comprender la propuesta de solución. Se parte de los conceptos 
generales de la Ingieniería de Software para luego enfocarse en una técnica específica, 
modelos característicos, para finalmente concluir la transición del conocimiento hacia el
dominio de la planificación y variabilidad educativa.

\subepigraph{Conceptos asociados al tema}{
    
    En este epígrafe se desglosan los conceptos esenciales que convergen en el
    núcleo del estudio: desde los fundamentos literarios y el valor pedagógico de la
    fábula, pasando por los desafíos y oportunidades de su digitalización, hasta los
    sistemas tecnológicos diseñados para su gestión y análisis. La clarificación de
    estos términos no solo delimita el campo de acción, sino que también
    proporciona el lenguaje común y la base teórica necesaria para el diseño de
    una solución coherente y efectiva.

    \subsubepigraph{Ingeniería de Líneas de Productos de Software (SPL)}{

        La Ingeniería de \ac{SPL} es una estrategia de
        ingeniería de software que busca producir familias de sistemas de software
        similares a partir de un conjunto común de activos centrales mediante un
        proceso gestionado de desarrollo y gestión de la variabilidad. El objetivo
        principal es lograr la reutilización a gran escala, lo que se traduce en una reducción
        de costos, menor tiempo de comercialización y mayor calidad de los
        productos \citep{sgbuzz_lps}. El concepto de
        \ac{SPL} surgió a finales de los años 90 como respuesta directa a la crisis de
        productividad que enfrentaba la industria del software a pesar de las técnicas
        de reutilización orientadas a objetos. La \ac{SPLE} formaliza una práctica que ya
        existía intuitivamente en muchas empresas: construir un software base
        (plataforma) para generar rápidamente múltiples productos similares con
        variaciones controladas \citep{academialab_linea_productos}

        \unnumberedsubsection{Conceptos Clave de SPL:}
        \begin{enumerate}
            \item \textbf{Línea de Productos:} Es la familia completa de sistemas de software
            que pueden ser generados. Representa el espacio de posibilidades
            definido por el conjunto de activos.
            \item \textbf{Activos Centrales (\textit{Core Assets}):} Son los componentes fundamentales
            (código, arquitectura, modelos de diseño, documentación, planes de
            prueba) diseñados específicamente para ser compartidos y reutilizados
            en toda la línea. La inversión inicial en estos activos es clave para el
            éxito de la estrategia.
            \item \textbf{Variabilidad:} Es la propiedad esencial de una \ac{SPL}. Define la capacidad
            de la línea para ser adaptada, personalizada o configurada para generar
            productos específicos que satisfagan las necesidades particulares de un
            cliente o escenario. La variabilidad es, en esencia, el conjunto de
            diferencias controladas entre los productos posibles.
            \item \textbf{Puntos de Variabilidad (\textit{Variability Points}):} Son las ubicaciones
            explícitas dentro de los Activos Centrales donde se deben tomar
            decisiones para configurar un producto. Estos puntos actúan como
            "interruptores" o "opciones" que guían el proceso de personalización.
        \end{enumerate}

        La motivación detrás de la \ac{SPL} es fundamentalmente económica y técnica,
        pues permite el ahorro de costos, ya que al compartir la gran mayoría de los
        activos (entre el 70\% y 90\% del sistema), se evita la reimplementación del
        mismo código y diseño para cada nuevo producto, además, provee de una reducción del tiempo de comercialización (\textit{Time-to-Market}) pues la creación de
        un nuevo producto se reduce de un ciclo completo de desarrollo a un simple
        proceso de configuración y ensamblaje de activos preexistentes y permite una
        mejora de la calidad, pues como los activos centrales, al ser reutilizados y
        probados en múltiples productos, alcanzan un nivel de madurez y robustez
        superior al que podría lograrse en un desarrollo de producto único \citep{ingenieria_software_lps}.

        La \ac{SPL} opera a través de dos procesos interdependientes que deben
        gestionarse de manera continua, siendo el primero de estos la Ingeniería del
        Dominio (\textit{Domain Engineering}), proceso hacia arriba que se enfoca en la
        creación y el mantenimiento de los Activos Centrales, cuya tarea principal es el
        modelado de la variabilidad para comprender y documentar todas las posibles
        configuraciones que la línea puede soportar y que constituye la base teórica y
        formal de la línea, y un segundo, la Ingeniería de la Aplicación (\textit{Application
        Engineering}), el cual es el proceso hacia abajo que utiliza los Activos Centrales
        para configurar y generar productos específicos, que implica tomar las
        decisiones de variabilidad definidas en el modelo (seleccionar las opciones)
        para producir una instancia concreta de la línea de productos. La herramienta
        principal y estándar de facto para formalizar y gestionar la variabilidad durante
        la Ingeniería del Dominio es el \ac{FM}.
        Este modelo establece el puente formal entre las opciones de la línea de
        productos y la capacidad de generar una instancia válida, lo cual es el principio
        que se busca trasladar al diseño curricular.
    }
    \subsubepigraph{Modelos de Características}{

        Es necesario contar con un modelo que represente de manera explícita las
        variaciones y opciones disponibles en la línea de productos. Este modelo
        debe ser capaz de capturar la complejidad de las relaciones entre las
        características y sus variaciones, así como las restricciones que pueden
        existir entre ellas \citep{pohl2005software}. La representación gráfica de estas relaciones es
        fundamental para facilitar la comprensión y el análisis de la variabilidad
        en la línea de productos.
    }
}}
